% Generated draft: requires mathematical and source review.
\documentclass[11pt]{article}
\usepackage[margin=25mm]{geometry}
\usepackage{fontspec}
\setmainfont{Latin Modern Roman}
\usepackage{amsmath,amssymb,mathtools,unicode-math}
\setmathfont{Latin Modern Math}
\usepackage{enumitem,needspace,xurl,hyperref}
\setlength{\parindent}{0pt}
\setlength{\parskip}{5pt}
\setlength{\emergencystretch}{3em}
\newcommand{\N}{\mathbb{N}}
\newcommand{\Z}{\mathbb{Z}}
\newcommand{\Q}{\mathbb{Q}}
\newcommand{\R}{\mathbb{R}}
\newcommand{\C}{\mathbb{C}}
\newcommand{\F}{\mathbb{F}}
\newcommand{\A}{\mathbb{A}}
\newcommand{\PP}{\mathbb{P}}
\newcommand{\E}{\mathbb{E}}
\newcommand{\eps}{\varepsilon}
\newcommand{\dd}{\,\mathrm{d}}
\begin{document}
\section*{UnsolvedMath ID: 30000791}
\textbf{Source problem code:} catalogue entry
\subsection{Definitions and mathematical statement}
Let \(S\) be a closed connected oriented surface of genus \(g\geq2\). Under the orientation-preserving convention, write
\[
 \operatorname{MCG}(S)=\operatorname{Homeo}^+(S)/\operatorname{Homeo}_0(S),
\]
where \(\operatorname{Homeo}_0(S)\) consists of homeomorphisms isotopic to the identity, and let \(q_H\) be the quotient homomorphism. For a subgroup \(\Gamma\leq\operatorname{MCG}(S)\), with inclusion \(\iota_\Gamma\), a homeomorphism lift is a group homomorphism
\[
 \rho_H:\Gamma\longrightarrow\operatorname{Homeo}^+(S),\qquad
 q_H\circ\rho_H=\iota_\Gamma.
\]
Equivalently, the restricted extension
\[
 1\longrightarrow\operatorname{Homeo}_0(S)
 \longrightarrow q_H^{-1}(\Gamma)
 \longrightarrow\Gamma\longrightarrow1
\]
splits. Such a lift is automatically injective. Choosing an individual representative for each mapping class is insufficient: representatives must satisfy all group multiplication relations exactly.

After fixing a smooth structure and a differentiability class \(r\), let
\(q_r:\operatorname{Diff}^{r,+}(S)\to\operatorname{MCG}(S)\) send a diffeomorphism to its topological isotopy class. A differentiable lift is a homomorphism
\[
 \rho_r:\Gamma\longrightarrow\operatorname{Diff}^{r,+}(S),\qquad
 q_r\circ\rho_r=\iota_\Gamma.
\]
The problem is to characterize all subgroups \(\Gamma\) admitting a homeomorphism lift, and to characterize those admitting a diffeomorphism lift in the intended regularity. These are distinct realization questions; the source gives no restriction to finite, finitely generated, or finite-index subgroups. The notation \(\operatorname{Diff}(S)\) does not specify \(r\); in the smooth convention it means \(r=\infty\).

The supplied assessment dated 2026-08-21 reports non-realizability of the whole mapping class group and realizability of many special subgroups, while the full subgroup classification remains open. Its qualitative summary does not supply the exact genus and regularity hypotheses of each cited theorem.
\subsection{Short English statement}
The question asks which collections of mapping classes can be represented simultaneously by an actual group action on the surface. It asks both for actions by homeomorphisms and for actions by diffeomorphisms, with regularity and orientation conventions requiring explicit agreement with the source.
\subsection{Sources}
\begin{itemize}
\item Local UnsolvedMath registry, permanent ID \(30000791\), exact source code \texttt{catalogue entry}, proposed year 2007; no individual proposer is supplied.
\item \emph{Dynamische Systeme (2008)}, Oberwolfach Reports, DOI 10.4171/owr/2007/33, pp. 1913--1956. \url{https://doi.org/10.4171/owr/2007/33}.
\item Supplied literature reference: Markovic--Saric, ``non-realizability,'' \url{https://arxiv.org/abs/0807.0182}.
\item The supplied \texttt{corrected\_verified} note reports checking mathematical content against the official report. The open label and assessment checked 2026-08-21 are dated supplied triage requiring independent verification of current status and exact cited hypotheses.
\end{itemize}
\end{document}
